\section{Stacks}

\solution{Problem}{32}
The solution is Listing~\ref{lst:balanced} of Chapter~\ref{ch:stacks}.
\idea Push openers; a closer must match the most recent unmatched opener, which is the
top. Fail on a mismatch, on a closer with an empty stack, or on openers left at the end.
\complexity $\Theta(n)$ time, $O(n)$ stack.
\mistakes Returning 1 at the end without checking that the stack is empty (\cc{"(("});
calling \cc{top()} without first checking \cc{empty()}, so that \cc{"))"} reads from an
empty stack; using counters, which accept \cc{"([)]"}.

\solution{Problem}{33}
The solution is Listing~\ref{lst:postfix} of Chapter~\ref{ch:stacks}; the table there
traces \cc{2*(3+4)-5}.
\idea Shunting-yard for the conversion; one operand stack for the evaluation, where the
\emph{first} pop is the right operand.
\complexity Both $\Theta(n)$: every token is pushed and popped at most once.
\mistakes Popping operators of equal precedence too late (\cc{>} instead of \cc{>=}: then
\cc{8-3-2} evaluates to 7); forgetting to discard the \cc{(}; swapping the operands of
\cc{-} and \cc{/} (Tracing~\ref{trace:13}). For the bonus: \cc{\^} has the highest
precedence and is right-associative, so an incoming \cc{\^} pops only operators of
\emph{strictly} higher precedence.

\solution{Problem}{34}
\idea Keep a second stack \cc{mins} in step with \cc{all}: its top is always the minimum
of the whole stack. Each push pushes $\min(x,\text{top of mins})$; each pop pops both.
\listingfile{sol/p34.cpp}{A min stack.}
\complexity Every operation $\Theta(1)$; $2n$ cells. In the bonus version \cc{mins} is
pushed only when $x\le$ its top and popped only when the popped item equals its top; the
test must be $\le$ so that a repeated minimum (pushes 3, 3, pop) is still on \cc{mins}
after one of its copies leaves.
\mistakes A single variable for the minimum (it cannot be restored after a pop); popping
\cc{all} without popping \cc{mins}.

\solution{Problem}{35}
The solution is Listing~\ref{lst:nextgreater} of Chapter~\ref{ch:stacks}, with
Theorem~\ref{thm:monostack} for the running time.
\idea The stack holds the indices still waiting for a greater value; \cc{a[i]} answers
all smaller waiting values, which are on top.
\complexity $\Theta(n)$: every index is pushed and popped at most once. Circular bonus:
run the loop for $i=0..2n-1$ on \cc{a[i \% n]} and push only in the first round.
\mistakes Storing values instead of indices (then the answer cannot be written to the
right cell); using \cc{<=}, which treats an equal value as greater (\cc{[1,1]} would give
\cc{[1,-1]}).

\solution{Problem}{36}
\idea The same monotonic stack: the waiting days, with non-increasing temperatures.
When day $i$ is strictly warmer than the day $j$ on top, day $j$ waited $i-j$ days.
\listingfile{sol/p36.cpp}{Daily temperatures.}
\complexity $\Theta(n)$ time.
\mistakes Using \cc{<=} (equal temperatures are not warmer); forgetting to write 0 for the
days left on the stack (here written explicitly, since \cc{wait} is not initialised).

\solution{Problem}{37}
\idea The online ``previous greater'' problem. Keep pairs (price, span) on two stacks in
step, prices strictly decreasing from bottom to top. A new price absorbs every pair with a
price $\le$ it, together with the days that pair already covered.
\listingfile{sol/p37.cpp}{Online stock span.}
\complexity $\Theta(1)$ amortised per call: every pair is pushed once and popped at most
once (Tracing~\ref{trace:15}).
\mistakes Storing every price and scanning backwards ($\Theta(n)$ per call, $\Theta(n^2)$
in total for decreasing prices); using \cc{<} (an equal price must be absorbed).

\solution{Problem}{38}
\idea The stack holds the survivors so far, left to right. Only a left-mover arriving at a
right-mover on top can collide; it destroys smaller right-movers one after another, dies
with an equal one, or dies against a larger one. Left-movers with nothing to their left
that moves right survive. At the end the stack, read from bottom to top, is the answer;
popping it fills \cc{out} from the right.
\listingfile{sol/p38.cpp}{Asteroid collision.}
\complexity $\Theta(n)$: each asteroid is pushed and popped at most once.
\mistakes Letting a right-mover collide with a left-mover on top (\cc{[-2,-1,1,2]}: they
move apart); stopping after one collision; returning the survivors in reverse order.

\solution{Problem}{39}
\idea At each \cc{[} remember, on two stacks, the repeat count and the position in
\cc{out} where the bracketed part starts. Characters are written straight to \cc{out}. At
\cc{]} the part \cc{out[start..len-1]} is already there once; append it $k-1$ more times.
Nested brackets work because the inner part is expanded before the outer \cc{]} copies it.
\listingfile{sol/p39.cpp}{Decode string.}
\complexity Linear in the length of the output (every output character is written once).
\mistakes Reading only one digit of the count (\cc{10[a]}); not resetting the count after
\cc{[}; copying the part $k$ times instead of $k-1$.

\solution{Problem}{40}
\idea The stack holds the letters that survived so far. A letter equal to the top cancels
with it; otherwise it is pushed. Cancellations can cascade (\cc{abba}) because after a pop
the new top is the next older survivor.
\listingfile{sol/p40.cpp}{Removing adjacent duplicates.}
\complexity $\Theta(n)$.
\mistakes Removing only one pass of pairs (\cc{abbaca} becomes \cc{aaca}, not \cc{ca});
copying the stack to \cc{out} in reverse order.

\solution{Problem}{41}
\idea Insertion sort between two stacks. Pop $x$ from \cc{s}; move back to \cc{s} every
item of \cc{r} larger than $x$; push $x$ on \cc{r}. Then \cc{r} is sorted with the
largest on top; pouring it into \cc{s} puts the smallest on top.
\listingfile{sol/p41.cpp}{Sorting a stack with one extra stack.}
\complexity $\Theta(n^2)$ in the worst case (each item may be moved back and forth $O(n)$
times), $\Theta(n)$ if \cc{s} is already ordered suitably.
\mistakes Comparing with \cc{>=}, which still sorts but moves equal items needlessly;
forgetting the final pour, which leaves the largest item on top.

\solution{Problem}{42}
\idea Simulate: push the items of \cc{pushed} in order, and after each push pop as long as
the top is the next wanted item. The sequence is valid iff every item was popped. Bonus:
an item must be pushed before it is popped, and once the top is the wanted item, popping
now is never worse than pushing more (which would only bury it), so the greedy rule
never misses a solution (Section~\ref{sec:stackperm}).
\listingfile{sol/p42.cpp}{Validate stack sequences.}
\complexity $\Theta(n)$: $n$ pushes and at most $n$ pops.
\mistakes Popping only once after each push (\cc{4 5 3 2 1} needs several pops at the
end); reading \cc{popped[j]} with $j=n$ (guarded here by the empty stack after the last
pop).
